\documentclass[UTF8,fontset=none,11pt,a4paper]{ctexart}
\usepackage[left=21mm,right=21mm,top=27mm,bottom=22mm,headheight=26pt,headsep=6mm]{geometry}
\usepackage{fontspec}
\setmainfont{texgyretermes-regular.otf}[BoldFont=texgyretermes-bold.otf,ItalicFont=texgyretermes-italic.otf]
\setCJKmainfont{Songti.ttc}[Path=/System/Library/Fonts/Supplemental/,FontIndex=6,AutoFakeBold=2]
\setCJKsansfont{Songti.ttc}[Path=/System/Library/Fonts/Supplemental/,FontIndex=6,AutoFakeBold=2]
\usepackage[table]{xcolor}
\definecolor{TableHead}{HTML}{EEF0ED}
\definecolor{Muted}{HTML}{666C68}
\usepackage{booktabs,longtable,array,fancyhdr,needspace,xurl}
\usepackage[unicode,hidelinks]{hyperref}
\urlstyle{same}
\linespread{1.32}
\setlength{\parindent}{2em}
\setlength{\parskip}{3pt}
\setlength{\LTpre}{8pt}\setlength{\LTpost}{10pt}
\ctexset{section={format=\fontsize{15}{20}\selectfont\bfseries,beforeskip=18pt,afterskip=8pt},subsection={format=\fontsize{12}{16}\selectfont\bfseries,beforeskip=12pt,afterskip=6pt},subsubsection={format=\normalsize\bfseries,beforeskip=9pt,afterskip=4pt}}
\pagestyle{fancy}\fancyhf{}
\renewcommand{\headrulewidth}{.3pt}
\renewcommand{\footrulewidth}{0pt}
\fancyhead[L]{\parbox[b]{\linewidth}{\fontsize{8}{10}\selectfont\color{Muted}现金账户：银行存款还是政府货币基金·本次账户资料及办理记录}}
\fancyfoot[C]{\small\color{Muted}\thepage}
\hypersetup{pdftitle={本次账户资料及办理记录},pdfsubject={仅供模拟使用·现金账户：银行存款还是政府货币基金·20-03},pdfauthor={金融决策模拟材料}}
\begin{document}
{\noindent\fontsize{20}{25}\selectfont\bfseries 本次账户资料及办理记录\par}
\vspace{7pt}{\noindent\small\color{Muted}仅供模拟使用·现金账户：银行存款还是政府货币基金·20-03\par}\vspace{10pt}
\subsection*{A：银行存款归集}
\textbf{转入资金的去向：} 平台按存款归集规则转入参与银行，形成银行存款


\textbf{开始比较使用的收益指标：} 存款年化收益率APY为1.84\%


\textbf{指标的含义：} 按所列存款利率及复利规则计算的年化收益率，存款利率可调整


\textbf{基金内部运营费：} 不购买该基金，无该项基金费


\textbf{资金保障机制：} 符合记录及存款归集条件时，按参与银行的FDIC存款保险规则保障


\textbf{资金运作：} 参与银行依存款规则持有及计息


\textbf{价值及账户付款：} 从银行存款中提取相应金额处理账户付款


\textbf{使用账户的付款工具：} 借记卡、支票、账单支付、转账


\subsection*{B：政府货币市场基金}
\textbf{转入资金的去向：} 自动买入政府货币市场基金份额


\textbf{开始比较使用的收益指标：} 基金七日年化收益率为3.28\%


\textbf{指标的含义：} 根据最近7日净收益折算的年化指标，已反映基金运营费，随实际基金收益变化


\textbf{基金内部运营费：} 总年度运营费率0.42\%，从基金资产承担，不另外按账户余额开账单


\textbf{资金保障机制：} 基金份额不是银行存款，不属于FDIC存款保险


\textbf{资金运作：} 主要投资现金、美国政府证券及符合规则的回购协议


\textbf{价值及账户付款：} 基金以每份1美元的稳定价格为目标，账户付款时自动赎回相应份额；该价格目标不是本金保证


\textbf{使用账户的付款工具：} 相同，使用自动赎回功能处理付款


这里的两个收益指标是本次场景开始时采用的起始资料，不是固定12个月的合同利率。B列示的七日年化指标与A的APY采用不同计算口径；之后的利率、基金实际收益和余额随各自机制及你的支出变化。


\textbf{公开历史资料。} 你查阅的现金产品资料附有以下两项历史记录。它们涉及的机构和产品均不是本次账户的参与银行或政府货币基金：


\subsection*{资产及事件}
\textbf{2008年9月：} 货币基金R持有企业发行的债务证券，其中包括雷曼证券；基金每份净值降至0.97美元


\textbf{2023年3月：} 银行V倒闭，存款账户由接管及处置安排处理


\subsection*{当时的处理记录}
\textbf{2008年9月：} 基金宣布暂停支付赎回款，随后进入相关处置安排


\textbf{2023年3月：} 主管部门采用系统性风险例外保护该行全部存款，存款人自3月13日起可以使用资金；该事件包括超过常规保险限额的存款


\subsection*{办理确认事项}
A的FDIC保险按每个参与银行、每位存款人和账户所有权类别计算，通常每类最高250,000美元。该银行其他同类别存款需合并计算；本次已核对合并余额在额度内，平台保留满足归集保障条件的记录。保险对应参与银行倒闭时的合格存款，不是给基金投资提供的保障。


B采用政府货币市场基金规则，至少99.5\%的总资产配置于现金、政府证券或以现金／政府证券充分抵押的回购协议。基金收入随短期资产收益变化，分配的收益留在账户并按核心头寸机制继续运作。提款所用份额按照实际适用的赎回价值结算。


两项日常支付工具相同，本次没有因选择某项而增加专属客服或理财服务。你选择的是现金的持有和收益机制；以后可以按平台规则更换核心头寸，但本次不把将来换项安排预定为结果。

\end{document}